\documentclass[11pt]{article}

\usepackage[utf8]{inputenc}
\usepackage[T1]{fontenc}
\usepackage{geometry}
\geometry{a4paper, margin=1.05in}
\usepackage{hyperref}
\hypersetup{
  colorlinks=true,
  linkcolor=blue,
  citecolor=blue,
  urlcolor=blue,
  breaklinks=true
}
\usepackage{times}
\usepackage{natbib}
\usepackage{booktabs}
\usepackage{listings}
\usepackage{xcolor}
\usepackage{enumitem}
\usepackage{amsmath}
\usepackage{amssymb}
\usepackage{microtype}
\usepackage{xurl}

\lstset{
  basicstyle=\ttfamily\small,
  frame=single,
  breaklines=true,
  postbreak=\mbox{$\hookrightarrow$\space},
  backgroundcolor=\color{gray!10},
  showstringspaces=false
}

\title{Fine-Tuning Llama 3.2 for U.S. Immigration Law Q\&A Using AWS SageMaker}
\author{Nazarii Shportun \\
\texttt{nshportun@gmail.com}}
\date{May 2026}

\begin{document}
  \maketitle

  \begin{abstract}
    This paper describes an end-to-end pipeline for constructing a large-scale, source-grounded question-answering dataset covering U.S. immigration law, and for fine-tuning a small language model on that dataset. Starting from official U.S. government sources -- including the USCIS Policy Manual, federal regulations (8 CFR), BIA precedent decisions, and immigration statistics -- I crawl, parse, normalize, and chunk 10,056 canonical documents. Using Amazon Bedrock (Claude Sonnet 4.6), I generate 17,058 structured Q\&A pairs across 13 immigration subdomains, each annotated with source provenance, authority level, answer type, and immigration subtopic. I then fine-tune Meta's Llama 3.2 3B Instruct model via AWS SageMaker JumpStart using parameter-efficient LoRA, merge the adapters into the base weights, and publish both the dataset and model to Hugging Face. The complete pipeline -- from first crawl to published model -- runs end-to-end on commodity cloud infrastructure for approximately \$25 in total compute cost. All artifacts are publicly available at \url{https://github.com/natebell510/usa-immigration}.
  \end{abstract}

  \section{Introduction}
  \label{sec:intro}

  U.S. immigration law is one of the most consequential legal domains for millions of people each year. Whether navigating visa applications, naturalization procedures, asylum claims, or removal proceedings, individuals often face these processes without legal representation. Large language models have the potential to make this kind of information more accessible, but general-purpose models struggle with domain-specific procedural questions. They get form numbers wrong, misstate filing deadlines, and confuse similar-sounding immigration categories.

  The core issue is one of representation: immigration law is well documented in official government sources, but it is underrepresented in the training data used for most LLMs. This paper describes the construction of a new dataset, available as \url{https://huggingface.co/datasets/nshportun/usa-immigration-law-qa} on Hugging Face, containing 17,058 source-grounded Q\&A pairs. It also describes the fine-tuning of a Llama 3.2 3B model, available as \url{https://huggingface.co/nshportun/usa-immigration-llama-3.2-3b-v3}.

  \section{Related Work}
  \label{sec:related_work}

  I review three lines of related work: immigration legal datasets, legal domain fine-tuning, and parameter-efficient adaptation methods.

  \paragraph{Immigration legal datasets.}
  Publicly available immigration Q\&A datasets remain sparse. The \texttt{harshitha008/US-immigration-laws} collection on Hugging Face (Apache 2.0) provides approximately 8,900 QA pairs derived from scraped immigration law pages, but lacks source URLs, authority-level annotations, or subdomain labels. Law StackExchange hosts community Q\&A under an immigration tag, but the content is noisy and largely unsourced.

  \paragraph{Legal QA and fine-tuning.}
  LegalBench~\citep{guha2023legalbench} benchmarks LLMs on legal reasoning tasks but does not include a U.S. immigration-specific split. Prior work on legal domain fine-tuning (e.g., LexGPT, LawInstruct) has focused primarily on case law and contract analysis, not immigration procedure.

  \paragraph{Small model fine-tuning.}
  LoRA~\citep{hu2022lora} enables parameter-efficient fine-tuning by injecting trainable low-rank decomposition matrices into attention layers. This approach has become the dominant paradigm for domain-specific adaptation when compute budgets are tight.

  \section{Data Pipeline}
  \label{sec:data_pipeline}

  \subsection{Source Identification}
  I identified seven source tiers ordered by authority level (Table~\ref{tab:sources}).

  \begin{table}[ht]
    \centering
    \caption{Source tiers by authority level.}
    \label{tab:sources}
    \begin{tabular}{lllc}
      \toprule
      Tier & Source & Authority & Records \\
      \midrule
      1 & USCIS Policy Manual & primary\_official & $\sim$3,200 pages \\
      2 & USCIS Forms \& Instructions & primary\_official & 87 forms \\
      3 & 8 CFR / INA & primary\_official & Title 8, parts 1--1400 \\
      4 & BIA Precedent Decisions & primary\_official & Selected precedents \\
      5 & DHS/CBP Yearbook Statistics & primary\_official & Annual tables \\
      6 & harshitha008/US-immigration-laws & secondary\_reputable & 8,897 QA pairs \\
      7 & Law StackExchange & community & Curated threads \\
      \bottomrule
    \end{tabular}
  \end{table}

  \subsection{Crawling}
  Each source tier has a dedicated crawler in \texttt{scripts/crawl/}. All crawlers respect \texttt{robots.txt}, enforce a configurable delay (\texttt{CRAWL\_DELAY\_SECONDS=1.5}), and identify themselves with a descriptive user-agent string.

  Raw content is stored in S3 under a structured prefix \texttt{v1/raw/\{source\}/\{doc\_id\}}.

  \subsection{Parsing and Normalization}
  Raw files pass through format-specific parsers that extract clean text, then through a normalization layer that assigns a canonical \texttt{doc\_id}, attaches metadata, deduplicates, and chunks documents into 512-token windows with 64-token overlap. The output is a 10,056-document corpus stored in \texttt{canonical\_corpus\_validated.jsonl}.

  \subsection{Q\&A Generation via Amazon Bedrock}
  I used Claude Sonnet 4.6 via Amazon Bedrock to generate structured Q\&A pairs from corpus chunks. Each generated pair includes question, answer, answer type, source span, time sensitivity flag, and topic tags.

  \subsection{Dataset Splitting}
  I perform a stratified split by immigration subtopic: 16,065 pairs (94.2\%) for training and 993 pairs (5.8\%) for evaluation.

  \section{Fine-Tuning}
  \label{sec:finetuning}

  \subsection{Model Selection}
  I selected Meta Llama 3.2 3B Instruct because it fits on a single \texttt{ml.g5.2xlarge} instance with int8 quantization and supports the required chat format.

  \subsection{SageMaker JumpStart Configuration}
  Final hyperparameters include LoRA (\texttt{r=8}, \texttt{alpha=32}), \texttt{per\_device\_train\_batch\_size=1}, \texttt{max\_input\_length=512}, and \texttt{merge\_weights=True}.

  \subsection{Total Pipeline Cost}
  \begin{table}[ht]
    \centering
    \caption{Breakdown of pipeline costs.}
    \label{tab:costs}
    \begin{tabular}{lr}
      \toprule
      Component & Cost \\
      \midrule
      Bedrock Claude QA generation ($\sim$17K pairs) & $\sim$\$18 \\
      SageMaker ml.g5.2xlarge training ($\sim$1.25 hrs) & $\sim$\$2 \\
      S3 storage (crawl + artifacts) & $\sim$\$1 \\
      \midrule
      \textbf{Total} & \textbf{$\sim$\$21} \\
      \bottomrule
    \end{tabular}
  \end{table}

  \section{Results and Discussion}
  \label{sec:results}

  \subsection{Published Artifacts}
  \begin{itemize}
    \item \textbf{Dataset:} \url{https://huggingface.co/datasets/nshportun/usa-immigration-law-qa}
    \item \textbf{Model:} \url{https://huggingface.co/nshportun/usa-immigration-llama-3.2-3b-v3}
  \end{itemize}

  \subsection{Limitations}
  The model may not reflect regulatory changes after the crawl date. It should not be used as a substitute for legal counsel.

  \section{Conclusion}
  \label{sec:conclusion}
  A high-quality, source-grounded dataset for a specialized legal domain can be constructed end-to-end with commodity cloud tools for under \$25.

  \appendix
  \section{Sample Q\&A Pairs}
  \label{app:samples}

  \paragraph{Family-based immigration -- form mode.}
  \textbf{Q:} What evidence must a U.S. citizen petitioner submit with Form I-130 to prove their relationship to a spouse born abroad?

  \textbf{A:} A U.S. citizen petitioner must submit: (1) evidence of U.S. citizenship; (2) the marriage certificate; (3) proof that any prior marriages were legally terminated; and (4) passport-style photos of both parties.

  \section*{References}
  \begin{thebibliography}{9}

    \bibitem[Guha et al.(2023)]{guha2023legalbench}
    Guha, N., Nyarko, J., Ho, D., R\'e, C., et al. (2023). LegalBench: A Collaboratively Built Benchmark for Measuring Legal Reasoning in Large Language Models. \emph{arXiv preprint arXiv:2308.11462}.

    \bibitem[Hu et al.(2022)]{hu2022lora}
    Hu, E. J., Shen, Y., Wallis, P., et al. (2022). LoRA: Low-Rank Adaptation of Large Language Models. \emph{International Conference on Learning Representations (ICLR)}.

    \bibitem[Meta(2024)]{meta2024llama32}
    Meta AI. (2024). Introducing Llama 3.2: Advancing Vision and Small Models. Available at \url{https://ai.meta.com/blog/llama-3-2-connect-2024-vision-edge-mobile-devices/}.

  \end{thebibliography}

\end{document}